% =============================================================================
%  正文 3：§6 平代环观点、§7 两边对照
% =============================================================================

\section{平代环观点}
\label{sec:adelic}

\subsection{}

模形式的概念可以通过两步推广：（a）推广为\textbf{自守形式}，（b）再推广为\textbf{自守表示}。
（$\GL(n)$ 的）自守表示，是 $\GL(n,\bbA)$ 的一个不可约表示，
出现在某个具有特定增长性、光滑性与等变性的复值函数空间（在 $\GL(n,\bbA)$ 上）里。
Langlands 与 Jacquet–Langlands [JL] 发现了如何给这样的表示配上 Hecke 特征值，使得下图交换：
\[
\begin{array}{ccc}
\text{模形式} & \longrightarrow & \text{自守形式}\\[2pt]
\big\downarrow & & \big\downarrow\\[2pt]
\text{Hecke 特征值} & \xleftarrow{\quad} & \text{自守表示}
\end{array}
\]

这些构造要用到\textbf{平代环}（adèle）
\[
\bbA=\bbA_{\bbQ}=\bbR\times\prod_p{}'\bbQ_p .
\]
它是实数与所有素数 $p$ 处的 $p$ 进数的\emph{限制直积}：
元素序列 $(a_{\infty},a_2,a_3,a_5,\dots)$（$a_p\in\bbQ_p$）属于 $\bbA$，
当且仅当存在 $N>0$ 使得对所有 $p\ge N$ 都有 $a_p\in\bbZ_p$。
对象 $\bbA$ 在逐坐标乘法下构成局部紧环，它包含对角嵌入的有理数 $\bbQ$ 作为离散子环。
（即每个 $q\in\bbQ$ 对应序列 $(q,q,q,\dots)$。）
我们称 $\bbR$ 为\textbf{无穷分支}，其余部分 $\bbA_{\mathrm f}=\prod'_p\bbQ_p$ 为\textbf{有限平代环}。

群 $\GL(2,\bbA)$ 是局部紧的，并且分解为限制直积
\[
\GL(2,\bbA)=\GL(2,\bbR)\times\prod_p{}'\GL(2,\bbQ_p),
\]
意思是序列 $g_{\infty},g_2,g_3,g_5,\dots$ 属于 $\GL(2,\bbA)$，当且仅当对充分大的所有 $p$ 有 $g_p\in\GL(2,\bbZ_p)$。
对角嵌入的子群 $\GL(2,\bbQ)$ 在 $\GL(2,\bbA)$ 中是离散的，于是可以取商 $\GL(2,\bbQ)\backslash\GL(2,\bbA)$。
模形式可以看成这个商上的函数。事实上有一个自然的（虽然并不完全显然的）等同
\[
\GL(2,\bbZ)\backslash\GL(2,\bbR)\;\cong\;\GL(2,\bbQ)\backslash\GL(2,\bbA)/\Kf,
\]
其中
\[
\Kf=\prod_p\GL(2,\bbZ_p)
\]
是有限平代群 $\GL(2,\bbA_{\mathrm f})$ 的紧开拓子群。这个等同诱导同构
\begin{equation*}
X=\sL(2,\bbZ)\backslash\he\;\cong\;\GL(2,\bbQ)Z(\bbA)\backslash\GL(2,\bbA)/\sO(2)\times\Kf .
\tag{6.1.1}
\end{equation*}
此外流形 $X$ 带有定义在有理数域 $\bbQ$ 上的复代数簇结构。
带这个代数结构，$X$ 称为 \textbf{Shimura 簇}\en{9}。

\begin{yizhu}[平代环为什么能把所有素数一次收进来]
用一句话说：\textbf{平代环是把「对所有素数 $p$ 同时成立」变成「在一个对象上成立」的工具}。
模形式本身只带一个复变量 $z$，看不出 $p$ 的痕迹；
但它的 Fourier 系数 $a_n$ 里已经含了全部素数，而 Hecke 特征值 $\lambda_p$ 正是按素数编号的。
(6.1.1) 做的事，是把「模形式 + 所有 $p$ 处的局部数据」打包成一个 $\GL(2,\bbA)$ 上的函数，
于是「对所有 $p$ 的特征值」就变成一个整体表示在不同局部分量上的信息。
\end{yizhu}

\subsection{平代环下的自守形式}

如 \S\ref{sec:modforms} 那样，每个（全纯）尖点模形式 $f$ 提升为函数
\begin{equation*}
\hat f:\GL(2,\bbQ)Z(\bbA)\backslash\GL(2,\bbA)\to\bbC .
\tag{6.2.1}
\end{equation*}
这样得到的函数满足特定的增长\en{10}、光滑\en{11}与等变\en{12}性质。
$f$ 的尖点性可以用 $\hat f$ 的某些积分\en{13}改写。
一般来说，$\GL(n,\bbA)$ 上的尖点自守形式，就是满足类似增长\en{14}、光滑\en{15}、等变\en{16}与尖点\en{17}条件的函数
\begin{equation*}
\varphi:\GL(n,\bbQ)Z(\bbA)\backslash\GL(n,\bbA)\to\bbC .
\tag{6.2.2}
\end{equation*}
把尖点自守形式构成的向量空间记作 $A_0=A_0(\GL(n,\bbA))$。

\subsection{自守表示}

本节的初等参考文献有 [K]、[G1] 与 [G2]，更完整的是 [Bu] \S3.2。
给定自守形式 $\varphi$（6.2.2），所有平移 $R_g(\varphi)$（$g$ 取遍 $\GL(n,\bbA)$）张成向量子空间
\begin{equation*}
L^2\bigl(\GL(n,\bbQ)Z(\bbA)\backslash\GL(n,\bbA)\bigr)
\tag{6.3.1}
\end{equation*}
的一个子空间，其闭包记作 $V_\varphi$。群 $\GL(n,\bbA)$ 作用在 $V_\varphi$ 上。
若自守形式 $\varphi$ 来自一个尖点模形式 Hecke 特征形式（$n=2$），则表示 $V_\varphi$ 不可约。
于是人们想定义：自守表示就是 (6.3.1) 中出现的任何不可约表示。
这个计划有两个问题。第一，对 (6.3.1) 的任何合理分解都会出现连续谱，
这会带来极大的困难，终究必须处理；不过眼下我们绕开它，只考虑满足尖点条件的 $L^2$ 可测函数构成的 Hilbert 空间
\begin{equation*}
L^2_0\bigl(\GL(n,\bbQ)Z(\bbA)\backslash\GL(n,\bbA)\bigr),
\tag{6.3.2}
\end{equation*}
这个空间分解为 $\GL(n,\bbA)$ 的不可约酉表示的离散和\en{18}。

(6.3.2) 的每个不可约成分 $\pi$ 分解为 $\GL(n,\bbR)$ 与 $\GL(n,\bbQ_p)$ 的不可约酉表示的乘积
$\pi=\pi_{\infty}\otimes\bigotimes_p\pi_p$，
而每个 $\pi_p$ 带有与群 $\GL(n,\bbQ_p)$ 相应的 Hecke 代数 $H_p$（见下文）相联系的特征值。
这些特征值是这样构造的。尖点自守形式空间
\begin{equation*}
A_0=A_0(\GL(n,\bbA))
\tag{6.3.3}
\end{equation*}
是 (6.3.2) 的子空间，它也同样分解为不可约表示的（Hilbert 空间直）和，
每个不可约成分称为一个\textbf{尖点自守表示}\en{19}。
此外，(6.3.2) 中出现的不可约表示与 (6.3.3) 中出现的不可约表示之间有一个自然的双射：
若
\[
\pi:\GL(n,\bbA)\to\GL(V)
\]
是 (6.3.2) 中出现的不可约酉表示，则 $K$-有限向量\en{20}构成的子空间 $V^{(K)}$
构成一个出现在 (6.3.3) 中的不可约尖点自守表示 $\pi'$。
而且这个自守表示 $\pi'$ 分解为限制张量积\en{21}
\[
\pi'=\pi'_{\infty}\otimes\bigotimes_p\pi'_p,
\]
其中 $\pi'_p$ 是 $\GL(n,\bbQ_p)$ 的不可约（容许\en{22}）表示，
而 $\pi'_{\infty}$ 是一个 $(\mathfrak g_{\infty},K_{\infty})$ 模（见注 20）。

接下来的计划，是给每个成分 $\pi'_p$ 配上 Hecke 特征值。
并非每个这样的不可约张量积都会出现在 (6.3.2) 中；
各个因子以何种方式出现在 (6.3.2) 的不可约成分里，据信反映了深刻的数论事实。

\begin{yizhu}[原文疑误：(6.3.2) 的分子]
原文 (6.3.2) 写作
\[
L^2_0\bigl(\GL(n,\bbA)Z(\bbA)\backslash\GL(n,\bbA)\bigr),
\]
而 (6.3.1) 与 (6.2.2) 都是
\[
\GL(n,\bbQ)Z(\bbA)\backslash\GL(n,\bbA).
\]
按上下文，(6.3.2) 的分子应为 $\GL(n,\bbQ)$；译文按正确形式书写。
\end{yizhu}

\subsection{自守表示的 Hecke 特征值}

本节固定一个有限素数 $p<\infty$，令 $G_p=\GL(n,\bbQ_p)$，$K_p=\GL(n,\bbZ_p)$。
Hecke 代数 $H_p$ 是双陪集
\begin{equation*}
K_pgK_p\in K_p\backslash G_p/K_p
\tag{6.4.1}
\end{equation*}
的所有有限形式线性组合构成的向量空间；
等价地，它是 $G_p$ 上光滑（即局部常值）、紧支、$K_p$-双不变的函数构成的向量空间。
它以双陪集 (6.4.1) 的特征函数为一组基。$H_p$ 中的乘法可以写成卷积\en{23}：
\[
(f*h)(z)=\int_{G_p}f(x)h(x^{-1}z)\,dx .
\]
Hecke 代数 $H_p$ 作用在 $K_p$-有限局部常值函数 $\varphi$ 上：
\[
(h\varphi)(z)=\int_{G_p}h(g)\varphi(zg)\,dg .
\]
固定一个不可约（容许）表示 $\pi_p:G_p\to\GL(V)$，它作为尖点自守表示 $\pi$ 的 $p$-成分出现。
若子群 $K_p$ 在 $V$ 上作用有非平凡的固定子空间 $V^{K_p}$，则称 $\pi_p$ 是\textbf{非分歧}的或\textbf{球}的，
（也说 $\pi$ 在 $p$ 处非分歧）。
此时 $K_p$ 固定向量构成的子空间是一维的，可以解释为 $p$ 处 Hecke 算子的特征空间。
事实上\en{24}，Hecke 代数 $H_p$ 通过某个 $\lambda:H_p\to\bbC$ 作用在这个一维空间上，
而它恰是代数同态，称为表示 $\pi_p$ 对应的 \textbf{Hecke 特征值}。

Satake 证明了\en{25}，这种同态的全体可以自然等同于非零复数对角 $n\times n$ 矩阵的等价类；
两个矩阵等价当且仅当差一个元素置换（即等同于 $T(\bbC)$，$T$ 是 $\GL(n)$ 中对角矩阵构成的环面）。
而且这个 \textbf{Satake 参数}决定了表示 $\pi_p$。
若 $\diag(t_1,t_2,\dots,t_n)$ 是 $n$ 个非零复数构成的对角矩阵，令 $\mu:T(\bbQ_p)\to\bbC^{\times}$ 为同态
\[
\mu\begin{pmatrix}x_1&&&\\&x_2&&\\&&\cdots&\\&&&x_n\end{pmatrix}
=\prod_{i=1}^{n}t_i^{\val(x_i)} .
\]
则 $\pi_p=\Ind_T^{G_p}(\mu)$ 就是把该特征诱导到 $G_p$ 得到的结果。细节见综述 [K]。

Satake 参数 $t=\diag(t_1,t_2,\dots,t_n)\in T(\bbC)$ 也可以等同于 $\GL(n,\bbC)$ 中的一个半单共轭类。
这个类中任何元素的特征值都是 $t_1,t_2,\dots,t_n$（可能被置换），
（相对于 $\GL(n,\bbC)$ 的标准表示）我们把它看作附着在 $\pi_p$ 上的、所要的「Hecke 特征值」\en{26}。
按通常方式构造 $\zeta$ 函数
\[
\zeta_p(y)=\frac{1}{\det(I-yt)} .
\]
若 $\pi$ 是 $\GL(n,\bbA)$ 的自守表示，则定义
\[
L(\pi,s)=\prod_p\zeta_p(p^{-s}),
\]
其中乘积取遍使 $\pi_p$ 非分歧的所有素数 $p$\en{27}。

\subsection{小结}

最终，朗兰兹猜想断言：应该存在一种方式，给 $\Gal(\Qbar/\bbQ)$ 的每个表示 $\rho$ 配上一个 $\GL(n,\bbA)$ 的自守表示 $\pi$，
使得所得的 $L$ 函数（差常数因子与 $s$ 的平移）相等：
\[
L(\rho,s)=L(\pi,s).
\]
换句话说，$\Fr_p$ 在表示 $\rho$ 上的特征值，应该等于 Hecke 代数 $H_p$ 在表示 $\pi$ 的 $p$-成分 $\pi_p$ 上的特征值。

\begin{keytakeaways}
这一节把第 5 节的古典语言整体翻译成现代语言。
翻译词典只有三行：模形式 $\to$ 自守形式 $\to$ 自守表示；
Fourier 系数 $\to$ Hecke 特征值 $\to$ Satake 参数；
「对所有 $p$ 同时成立」$\to$ 「一个 $\GL(n,\bbA)$ 上的函数」。
第 6.5 节的两行等式 $L(\rho,s)=L(\pi,s)$ 就是整个朗兰兹猜想。
\end{keytakeaways}

\section{两边对照}
\label{sec:compare}

\subsection{$n=2$ 的情形}

Eichler–Shimura 定理可以重新解释为 $n=2$ 时猜想的证明。
把 Galois 表示与模形式联系起来的做法有好几种。
Shimura 最初的描述 [Sh1] 给每个模形式 $f$ 配一条椭圆曲线 $E_f$，
然后证明 Galois 群作用在 $E_f$ 的挠点上，由此得到所要的 Galois 表示；
它可以解释为椭圆曲线平展上同调 $H^1_{\et}(E_f,\Ql)$ 上的 Galois 表示。
不过现代的做法是把这个 Galois 表示认同为 (6.1.1) 中所述 Shimura 簇
\[
X=\GL(2,\bbQ)Z(\bbA)\backslash\GL(2,\bbA)/\sO(2)\Kf
\]
的上同调 $H^1_{\et}(X,\Ql)$ 中的那个表示。
事实上 Eichler 与 Shimura 证明了：$X$ 上存在局部系数系 $E_k$，使得
\begin{equation*}
H^1_P(X,E_k;\bbQ)\otimes_{\bbQ}\bbC\;\cong\;S_k\oplus\overline{S_k},
\tag{7.1.1}
\end{equation*}
其中 $S_k$ 表示权为 $k$ 的尖点形式构成的（有限维）向量空间。
这里 $H^1_P(X,E_k)$ 指映射 $H^1_c(X,E_k)\to H^1(X,E_k)$ 的像，
而 $H^1_c$ 表示紧支上同调\en{28}。
在同构 (7.1.1) 下，左边 Hecke 对应的作用与右边 Hecke 算子的作用一致。
若尖点形式 $f$ 是 Hecke 特征形式，则其复共轭 $\bar f$ 也是所有 Hecke 算子的特征函数，且有相同的 Hecke 特征值。
于是 Hecke 代数的作用把左边的上同调群分解为若干二维子空间。
由此推出，平展上同调 $H^1_{\et,P}(X,E_k;\Ql)$ 在 Hecke 对应的作用下分解为若干二维子空间，
每个对应唯一一个尖点 Hecke 特征形式 $f$，从而对应一个自守表示。
这样就终于给出了 $\Gal(\Qbar/\bbQ)$ 的（某些）二维表示与 $\GL(2,\bbA)$ 的（某些）尖点自守表示之间的对应。
在此对应下，Galois 特征值确实对应于 Hecke 特征值——这是困难的部分。
为证明这一点，需要 Galois 作用与 Hecke 作用在平展上同调上的某种关系，
而 Eichler–Shimura 关系提供了它：粗略地说，对应 $T_p$ 同构于 $\sigma_p$ 与其转置之和（作为对应）。

\subsection{$n\ge 3$ 的情形}

上述论证如此有力，以至于人们要问：为什么它不能直接推广？
问题在于，对 $n\ge 3$，整个程序彻底失效，因为所需的 Shimura 簇不存在。
例如 $n=3$ 时，类似的对象
\[
\GL(3,\bbQ)Z(\bbA)\backslash\GL(3,\bbA)/\sO(3)\Kf
\]
根本不是代数簇，事实上它是一个（实）五维流形。这个问题是本领域最大的谜团之一。

\subsection{$\GL(1)$ 又如何？}

\begin{yizhu}[原文此处空缺（一）]
原文第 7.3 节只有这个标题，正文是空的。作者在开头已声明本文是 work in progress。
这一节本该讲的是：$n=1$ 时朗兰兹对应就是类域论（class field theory）——
$\Gal(\Qbar/\bbQ)^{\mathrm{ab}}$ 上的 Hecke 特征（即 Dirichlet 特征）与 $\GL(1,\bbA)$ 的自守表示（即 Hecke 特征标）一一对应，
两边 $L$ 函数相等。这也是整个猜想的原型。作者没有写，译文不代拟。
\end{yizhu}

\begin{keytakeaways}
这一节回答了「几何朗兰兹为什么要换个做法」这个动机问题：
对 $n=2$，模空间（Shimura 簇）恰好存在，于是能用上同调把两边拼起来；
对 $n\ge 3$，这个模空间不存在，走不通。
几何朗兰兹的做法是把这个「模空间」换成曲线的 Hecke 栈——
不再数椭圆曲线，而是数曲线的 $G$-丛；同调的对象从 Shimura 簇的平展上同调换成了 Hecke 栈上的层。
本库已译的 Gaitsgory 大纲、Scholze Bourbaki 报告、Sandvik 讲义，都是在这条替代路线上做文章。
\end{keytakeaways}
